\documentclass[../../../uem_phy_std_a.tex]{subfiles}

\begin{document}
    \begin{enumerate}
        \item $\mathrm{R}_2$の電流を$i$（右向き）とすると，$\mathrm{R}_3$の電流は$I-i$（右向き）である．%
            キルヒホッフ則より，
            
            \myspaceb%
            $\displaystyle%
                \Biggl\{
                    \begin{myaligna}
                        & 2Ri + 2RI = 3V_0 \; , \\
                        & 2Ri = R(I-i) + V_0
                    \end{myaligna}
                \qquad \therefore \; i = \myfracc{5V_0}{8R} \;,
                \quad I = \uwave{\myfracc{7V_0}{8R}} \;.
            $
        \item $\mathrm{R}_1$の電流を$i$（右向き），$\mathrm{R}_5$の電流を$j$（上向き）とすると，%
            $\mathrm{R}_2$の電流は$i+j$（右向き），$\mathrm{R}_3$の電流は$I-i$（右向き），%
            $\mathrm{R}_4$の電流は$I-i-j$（右向き）である．キルヒホッフ則より，
            
            \myspaceb%
            $\displaystyle%
                 \Biggl\{
                    \begin{myaligna}
                        & R(i+j) + Ri = 3R(I-i-j) + R(I-i) = V_0 \; , \\
                        & Ri = Rj + R(I-i)
                    \end{myaligna}
            $

            \myspaceb%
            $\displaystyle%
                \therefore \; i = \myfracc{4V_0}{9R} \;,
                \quad j = \myfracc{V_0}{9R} \;,
                \quad I = \uwave{\myfracc{7V_0}{9R}} \;.
            $        
    \end{enumerate}
\end{document}
